% s-183-p183-m002; original PDF 183
\Aphorism{s-183-p183-m002}{272}
Der Süden ist die Erlaubniss zur Buntheit, das hohe Lied der Chromatik. Ich sah bei Genua ein prangend rosarothes Haus mit tiefblauen Fensterläden, daneben eine Pinie — hätten wir Nordländer den Muth zu solchem schreienden, nicht tönenden Farbendreiklang, und wenn den Muth, hätten wir das Recht? Wär’s nicht unter unserer Sonne, die alle Farben Lügen straft, ein grenzenlos Erdreusten und ein zweckloses dazu? Aber hier ist das Unwahrscheinlichste gestattet; es giebt im Süden keinen schlechten Geschmack, während bei uns ausser Grau, Braun, Schwarz (den farblosen Farben) fast alles schlechter Geschmack ist. Im Süden putzen Weiber und Häuser sich excessiv, für die Excesse des vollen Lichts; wenn ein Nordländer von vollem Licht, von \Foreign{plein air} redet, so bewundern wir seine Bescheidenheit — ihn berauscht schon ein Fingerhut des Weines, der hier aus Fässern strömt!
\par

